\section{Stability for cycles}\label{sec:cycles}

The cycle construction gives another infinite family. Its additional
symmetry allows a short induction in place of the finite reduction
needed for paths.

\begin{theorem}\label{cycle:stable}
Let $C_j$ be a cycle with $j\ge3$ vertices and edges, or the graph with
two vertices and two parallel edges when $j=2$. Put $b_c=2$ at every
vertex. Then $X(C_j)$ is a smooth toric Fano variety with stable tangent
bundle, dimension $4j$, and Picard number $5j$.
\end{theorem}

\begin{proof}
The geometric assertions and connectedness of the ray matroid follow
from Proposition~\ref{graph:geometry}. Number vertices and edges modulo
$j$, orienting edge $i$ from $i$ to $i+1$, and write
$g_c=f^c_1+f^c_2$. The twisted base ray at $c$ is
\[
 f^c_0=-g_c+r_{c-1}-r_c.
\]
Rotation sends vertex $c$ to $c+1$ and hexagon $i$ to $i+1$, preserving
the coordinates in each transported summand. Reflection sends $c$ to
$-c$ and $i$ to $-i-1$, and negates both coordinates in each transported
hexagon. Both maps send the displayed twisted ray to the twisted ray
at the image vertex. They preserve adjacent pairs of hexagon rays and
the choices of an omitted base ray, hence preserve the fan. The same
reflection works for $j=2$, where it exchanges the parallel edges and
fixes the vertices.

Together with the independent hexagon reflections
$r_i\mapsto r_i$, $q_i\mapsto-r_i-q_i$ and the transpositions of
$f^c_1,f^c_2$, these maps give four ray orbits: the $2j$ distinguished
hexagon rays $\pm r_i$, the other $4j$ hexagon rays, the $2j$ untwisted
base rays, and the $j$ twisted base rays. A ray-spanned invariant
subspace containing any nondistinguished hexagon ray must contain the
whole plane of every hexagon. Thus it suffices to consider a common
hexagon choice $H\in\{0,1,2\}$, for empty, line, or plane, and common
binary choices $C,T$ for untwisted and twisted base rays.

Let $a,b,c$ be the normalized degrees of a distinguished hexagon ray,
a nondistinguished hexagon ray, and a base ray, respectively. Symmetry
gives the first two equalities across factors. Principal divisor
relations give equal degrees for all three base rays at a vertex, so
there is only one base value $c$. Their sum satisfies
\begin{equation}\label{cycle:weightsum}
 2a+4b+3c=4.
\end{equation}
The nine proper nonempty configurations are listed in
Table~\ref{cycle:flats}. For their ranks, selected untwisted base rays
contribute $2jC$, and hexagon choices contribute $jH$. If $C=0$, each
selected twisted ray has an independent base coordinate. If $C=T=1$,
its residual vector is $r_{c-1}-r_c$: these vectors span a space of
dimension $j-1$ when $H=0$, and vanish modulo the selected hexagon
lines when $H>0$. This gives every rank in the table. The other three
configurations are empty or have full rank.

\begin{table}[ht]
\centering
\caption{Proper cycle configurations and their weights.}
\label{cycle:flats}
\begin{adjustbox}{max width=\textwidth}
\begin{tabular}{ccccc}
\toprule
$H$ & $C$ & $T$ & Rank & Weight divided by $j$\\
\midrule
0&0&1&$j$&$c$\\
0&1&0&$2j$&$2c$\\
0&1&1&$3j-1$&$3c$\\
1&0&0&$j$&$2a$\\
1&0&1&$2j$&$2a+c$\\
1&1&0&$3j$&$2a+2c$\\
1&1&1&$3j$&$2a+3c$\\
2&0&0&$2j$&$2a+4b$\\
2&0&1&$3j$&$2a+4b+c$\\
\bottomrule
\end{tabular}
\end{adjustbox}
\end{table}

We will prove the bounds
\begin{equation}\label{cycle:bounds}
 a<\frac12,\qquad b>\frac14,\qquad \frac23<c<\frac56.
\end{equation}
They imply every strict weight inequality in the table. For example,
$3jc<(5/2)j\le3j-1$ for $j\ge2$; the row $(1,1,1)$ follows from
$2a+3c=4-4b<3$; and the rows with $H=2$ follow from
$2a+4b=4-3c$. The remaining inequalities follow directly from
\eqref{cycle:bounds}. Proposition~\ref{prop:symmetry} will then prove
stability.

Let $W,\kappa$ be the three-dimensional matrices in
\eqref{path:matrices}, and put $M=W\kappa$. For $0\le p,q\le2$, set
\[
 (W_r)_{pq}=(-1)^{p+q},\qquad
 (W_-)_{pq}=\frac{(-1)^{p+q}}{p+q+1},\qquad
 \kappa_1=\begin{pmatrix}3&-1&0\\1&0&0\\0&0&0\end{pmatrix}.
\]
Expanding the cyclic product of kernels and integrating each variable
contracts their coefficients cyclically. Consequently the volume and
facet integrals of Section~\ref{sec:graphs} are
\begin{align*}
 Z_j&=\operatorname{tr}(M^j),&
 A_j&=\operatorname{tr}(W_r\kappa M^{j-1}),\\
 B_j&=\operatorname{tr}(W_-\kappa M^{j-1}),&
 C_j'&=\operatorname{tr}(W\kappa_1M^{j-1}).
\end{align*}
Here $a=A_j/Z_j$, $b=B_j/Z_j$, $c=C_j'/Z_j$. The denominator is the
integral of the strictly positive cyclic kernel, so $Z_j>0$.
Define four rational sequences by
\[
 s_j^{(1)}=\tfrac12Z_j-A_j,\quad
 s_j^{(2)}=B_j-\tfrac14Z_j,\quad
 s_j^{(3)}=C_j'-\tfrac23Z_j,\quad
 s_j^{(4)}=\tfrac56Z_j-C_j'.
\]
The characteristic polynomial of $M$ is
\[
 x^3-\frac{27}{2}x^2+\frac{2491}{240}x-\frac{127}{864}.
\]
By Cayley--Hamilton, each sequence satisfies
\begin{equation}\label{cycle:recurrence}
 s_{j+3}=\frac{27}{2}s_{j+2}-\frac{2491}{240}s_{j+1}
             +\frac{127}{864}s_j.
\end{equation}
The starting values, obtained from the displayed matrices, are in
Table~\ref{cycle:seeds}. Index one denotes a formal trace value, not a
looped graph construction.

\begin{table}[ht]
\centering
\caption{Exact starting values for the four cycle inequalities.}
\label{cycle:seeds}
\begin{adjustbox}{max width=\textwidth}
\renewcommand{\arraystretch}{1.55}
\begin{tabular}{ccccc}
\toprule
$j$&$s_j^{(1)}$&$s_j^{(2)}$&$s_j^{(3)}$&$s_j^{(4)}$\\
\midrule
\csname @@input\endcsname tables/cycle_seeds.tex
\bottomrule
\end{tabular}
\end{adjustbox}
\end{table}

All first terms are nonnegative, all second and third terms are positive,
and componentwise
\[
 s_3-12s_2=
 \left(\frac{20887}{720},\frac{21871}{2880},
       \frac{19903}{2160},\frac{34037}{4320}\right)>0.
\]
Suppose $s_k>12s_{k-1}>0$ and $s_{k-2}\ge0$. Recurrence
\eqref{cycle:recurrence} gives
\[
 s_{k+1}>
 \left(\frac{27}{2}-\frac{2491}{2880}\right)s_k>12s_k>0,
\]
since $27/2-2491/2880-12=1829/2880>0$. Induction proves positivity
for each sequence at every $j\ge2$. These four signs are exactly
\eqref{cycle:bounds}, completing the proof.
\end{proof}

All data used by the induction are printed above. Exact checks reproduce
the seeds and compare the trace degrees with the independently integrated
graph formulas for $j=2,3$.\footnote{The checks are implemented in
\path{figures/code/compute06_cycle_tables.py}.}
These checks verify the finite matrix calculations; the induction proves
the assertion for every length.
